\documentclass[conference]{IEEEtran}
\IEEEoverridecommandlockouts

\usepackage{cite}
\usepackage{amsmath,amssymb,amsfonts}
\usepackage{graphicx}
\usepackage{textcomp}
\usepackage{booktabs}
\usepackage{algorithm}
\usepackage{algpseudocode}
\usepackage{url}
\usepackage{hyperref}

% ---------------------------------------------------------------
% PAGE BUDGET (6-page hard limit, two-column IEEEtran conference
% format is roughly 900-950 words per column, ~1900 words/page).
% Use this as a planning guide while drafting; delete once final.
%   I.    Introduction ................ 0.75 page
%   II.   Related Work ................ 0.75 page
%   III.  Problem Formulation ......... 1.00 page
%   IV.   Methodology (RL + GNN) ...... 1.25 page
%   V.    Baselines ................... 0.50 page
%   VI.   Experimental Setup .......... 0.75 page
%   VII.  Results & Discussion ........ 0.75 page  (+ 1 table, 1-2 figs)
%   VIII. Conclusion ................... 0.25 page
%   References .......................  fits in remaining space
%   TOTAL .............................. ~6.0 pages
% ---------------------------------------------------------------

\begin{document}

\title{Reinforcement Learning for Grid-Aware Electric Vehicle\\
Charging Station Placement in Urban Networks\\
\large[Working Title -- Replace With Final Title]}

% --- Author block -------------------------------------------------
% IEEE conference style hides author identity blocks are fine for
% camera-ready but check the target venue's double-blind policy
% before submission (many RL/power-systems venues are single-blind).
\author{
\IEEEauthorblockN{First Author\IEEEauthorrefmark{1}, Second Author\IEEEauthorrefmark{1}}
\IEEEauthorblockA{\IEEEauthorrefmark{1}Department/Institution, City, Country\\
Email: \{author1, author2\}@institution.edu}
}

\maketitle

% ===================================================================
\begin{abstract}
% TODO (~150-200 words): 1 sentence problem statement (EV charging
% station siting under budget + power-grid capacity constraints);
% 1-2 sentences on the proposed approach (MDP formulation, GNN-based
% policy over the street/grid graph, DRL training); 1 sentence on
% baselines compared against (genetic algorithm, greedy-benefit,
% greedy-demand); 1-2 sentences on the evaluation setting (Hanoi
% districts: Cau Giay, Dong Da, Nam Tu Liem, Tay Ho -- real road
% network + synthesized power grid); 1 sentence quantifying the
% headline result (e.g. X\% reduction in waiting time / Y\% higher
% coverage at equal budget vs. the strongest baseline).
\end{abstract}

\begin{IEEEkeywords}
electric vehicle charging, facility location, reinforcement learning,
graph neural networks, power distribution networks, urban infrastructure planning
\end{IEEEkeywords}

% ===================================================================
\section{Introduction}
% TODO (~0.75 page):
% - Motivate EV adoption growth and the charging-infrastructure siting
%   problem as a bottleneck.
% - State why siting is NOT just a coverage/demand problem: new
%   stations draw load from the existing distribution grid, and
%   substations/feeders have finite spare capacity -- placement must
%   be grid-aware or it creates overload risk.
% - Position against prior siting work (typically ignores grid
%   constraints) and prior grid-aware work (typically uses static
%   optimization, not sequential/adaptive decision-making).
% - State contributions as a bullet list, e.g.:
%   1) An MDP formulation of sequential station placement jointly
%      over spatial demand and grid capacity constraints.
%   2) A GNN-based policy/value architecture that consumes the road
%      network graph directly (node features, edge index, edge attr).
%   3) A city-scale power grid simulator (feeders, substations, POI
%      load) coupled to the placement environment as a hard/soft
%      constraint source.
%   4) An empirical comparison against genetic-algorithm and greedy
%      baselines across four real Hanoi districts.

% ===================================================================
\section{Related Work}
% TODO (~0.75 page), suggested subsections or a single flowing
% section with 3 paragraphs:
% A. EV charging station placement (facility location, MILP/greedy
%    heuristics, GIS-based siting studies).
% B. RL / metaheuristics for infrastructure placement (GA, PSO,
%    and DRL approaches to siting/routing problems).
% C. Grid-aware placement / hosting-capacity analysis (power-flow
%    constrained siting, distribution network impact studies).
% End with 1-2 sentences on the gap this paper fills.

% ===================================================================
\section{Problem Formulation}
\label{sec:problem}
% TODO (~1.0 page)

\subsection{System Model}
% - City road/POI graph $G=(V,E)$: nodes = intersections/POIs with
%   features (land price, demand, existing infrastructure); edges =
%   road segments with distances.
% - Power grid model: substations/feeders/buses with available
%   capacity (MW); mapping from graph node -> nearest bus.
% - Existing charging plan (pre-placed stations) as the environment's
%   starting state; budget $B$ as a depleting resource.

\subsection{MDP Formulation}
% Define formally:
% - State $s_t$: current plan (station locations + config), remaining
%   budget, per-node demand/coverage features, grid headroom.
% - Action $a_t$: choose a node to place/upgrade a station and a
%   charger configuration (or "stop").
% - Transition: budget and grid capacity update deterministically
%   given the chosen action.
% - Reward $r_t$: state the actual reward decomposition used in
%   \texttt{StationPlacementEnv.py} (score delta term + best-score
%   bonus term -- NOTE: verify in code before publishing, see \S
%   reward-audit item below) plus any grid-violation penalty term.
% - Episode termination: budget exhausted or max steps reached.

\begin{equation}
\label{eq:reward}
r_t = \Delta \text{score}_t \;+\; \lambda \cdot \mathbb{1}[\text{grid violation}] \quad \text{(fill in exact form)}
\end{equation}

% ===================================================================
\section{Methodology}
% TODO (~1.25 page)

\subsection{Graph-Based Observation Encoding}
% Describe graph\_features.py's node/edge feature construction and
% the normalization/scaling pipeline (preprocess_scaling.py /
% scaling\_constants.json) -- state what is normalized and why
% (comparability across districts of different size/demand).

\subsection{GNN Feature Extractor}
% Describe gnn\_extractor.py architecture: input node features,
% adjacency construction from edge\_index/edge\_attr, graph
% convolution layers, pooling into a fixed-size embedding, and how
% it is fused with global (budget, step count) features before the
% policy/value heads.

\subsection{RL Algorithm and Training}
% Name the DRL algorithm (e.g., DQN, given target\_update\_interval
% appears in configs) and give: network update rule, exploration
% schedule, replay buffer, key hyperparameters (learning rate,
% gamma, batch size, target update interval, max\_grad\_norm),
% and training compute budget (steps, wall-clock, hardware).

\begin{algorithm}[t]
\caption{Grid-Aware Station Placement (Rollout)}
\label{alg:rollout}
\begin{algorithmic}[1]
\State Initialize plan from existing infrastructure, budget $B$
\While{budget remains and not terminal}
    \State Observe graph state $s_t$ (node feats, edges, global feats)
    \State $a_t \leftarrow \pi_\theta(s_t)$ \Comment{GNN policy}
    \State Apply $a_t$: place/upgrade station, deduct cost
    \State Update grid headroom at affected bus
    \State Compute reward $r_t$ (Eq.~\ref{eq:reward})
\EndWhile
\end{algorithmic}
\end{algorithm}

% ===================================================================
\section{Baselines}
% TODO (~0.5 page): briefly describe, one short paragraph each:
% - Greedy-Benefit (greedy_benefit.py): iteratively picks the
%   node/config maximizing marginal social-benefit per cost.
% - Greedy-Demand (greedy_demand.py): iteratively picks the
%   highest-unmet-demand node.
% - Genetic Algorithm (algorithms/ga/train_ga.py): chromosome
%   encoding of the plan, fitness = norm\_score, selection/
%   crossover/mutation operators, population size, generations.
% State explicitly that all methods are evaluated on identical
% start-of-episode conditions (same existing plan file, same graph,
% same budget, same seed(s)) -- confirm this in code before writing
% it (see audit note below).

% ===================================================================
\section{Experimental Setup}
% TODO (~0.75 page)
\subsection{Datasets}
% Four Hanoi districts: Cau Giay, Dong Da, Nam Tu Liem, Tay Ho.
% Table with per-district stats: \#nodes, \#existing stations,
% \#substations/buses, budget.

\subsection{Power Grid Simulation}
% One paragraph: how the synthetic grid (feeders, substations, POI
% load) was generated / calibrated, and what "grid violation" means
% operationally (bus overload threshold, feeder capacity, etc).

\subsection{Metrics}
% Define each metric precisely, matching the actual code:
% - Coverage: fraction of demand nodes within service radius.
% - Waiting time: mean vs. max over stations -- STATE WHICH ONE,
%   the current code names variables inconsistently with a
%   ``mean'' docstring while computing a maximum; pick the true
%   statistic and document it accurately here.
% - Travel time, social benefit / fairness term, budget utilization.

\subsection{Reproducibility}
% Seeds used per run, number of independent seeds/replications
% (recommend $\geq$5 if only 1 is currently used), hardware/software
% versions, and a link/footnote to the code repository.

% ===================================================================
\section{Results and Discussion}
% TODO (~0.75 page + 1 table + 1-2 figures)

\begin{table}[t]
\centering
\caption{Performance comparison across districts (fill in; report mean $\pm$ std over $\geq$5 seeds).}
\label{tab:results}
\begin{tabular}{lcccc}
\toprule
Method & Coverage & Wait (min) & Travel (min) & Cost util. \\
\midrule
Greedy-Demand   & -- & -- & -- & -- \\
Greedy-Benefit  & -- & -- & -- & -- \\
Genetic Alg.    & -- & -- & -- & -- \\
\textbf{RL (ours)} & -- & -- & -- & -- \\
\bottomrule
\end{tabular}
\end{table}

% Figure placeholder: e.g. learning curve, or map of chosen station
% locations overlaid on the district graph (visualise.py output).
\begin{figure}[t]
\centering
% \includegraphics[width=\columnwidth]{figures/learning_curve.pdf}
\caption{Placeholder -- e.g., training reward curve or a map of the
resulting station placement for one district.}
\label{fig:placeholder}
\end{figure}

% Discussion: where RL wins/loses vs. GA/greedy, and WHY (sample
% efficiency vs. solution quality trade-off, sensitivity to grid
% constraints, generalization across districts of different size).

% ===================================================================
\section{Conclusion and Future Work}
% TODO (~0.25 page): 2-3 sentence summary of contribution + result,
% 2-3 sentences on limitations (single-city, synthetic grid, fixed
% budget model) and future work (multi-objective, real grid data,
% transfer across cities).

% ===================================================================
\section*{Acknowledgment}
% Optional -- funding, data providers, etc.

\bibliographystyle{IEEEtran}
\bibliography{references}

\end{document}
